%
% teil1.tex -- Beispiel-File für das Paper
%
% (c) 2020 Prof Dr Andreas Müller, Hochschule Rapperswil
%
% !TEX root = ../../buch.tex
% !TEX encoding = UTF-8
%
\subsection{Einfache Herleitung für Kreisbahnen
    \label{kepler:subsection:kreisbahnen}}

Ein besonders anschaulicher Weg zur Herleitung ergibt sich unter der vereinfachenden Annahme einer perfekt kreisförmigen Umlaufbahn \cite{astro2012}.
Auch wenn dies ein Spezialfall der Ellipse mit einer Exzentrizität von null ist, liefert dieser Ansatz bereits das korrekte Resultat \cite{astro2012}.
Ein Planet der Masse \(m\) bewegt sich dabei mit konstanter Geschwindigkeit auf einer Kreisbahn mit Radius \(r\) um eine deutlich schwerere zentrale Masse \(M\) \cite{astro2012}.
Obwohl der Betrag der Geschwindigkeit konstant bleibt, ändert sich fortwährend ihre Richtung, was eine ständige Beschleunigung zum Zentrum hin bedeutet \cite{astro2012}.
Für diese \emph{Zentripetalbeschleunigung} gilt der Zusammenhang
\begin{equation}
    a_c = \omega^2 r = \left(\frac{2\pi}{T}\right)^2 r ,
    \label{kepler:equation:zentripetal}
\end{equation}
wobei \(a_c\) die Beschleunigung und \(\omega\) die Winkelgeschwindigkeit bezeichnen \cite{astro2012}.
Da die Newtonsche Gravitationskraft genau diese Beschleunigung verursacht, lässt sich das Kraftgleichgewicht aufstellen und man erhält
\begin{equation}
    \frac{GM}{r^2} = \frac{4\pi^2}{T^2} r .
    \label{kepler:equation:gleichgewicht}
\end{equation}
Durch einfaches Umstellen dieser Formel ergibt sich direkt das dritte keplersche Gesetz für den Radius \(r\), welcher bei einer Kreisbahn exakt der grossen Halbachse \(a\) der Ellipse entspricht \cite{astro2012}.


\subsection{Allgemeine Herleitung für elliptische Bahnen
    \label{kepler:subsection:ellipsen}}

Für eine exakte Herleitung bei tatsächlichen, elliptischen Planetenbahnen startet man am besten mit dem Flächensatz, also dem zweiten keplerschen Gesetz \cite{astro2012}.
Die Ableitung der überstrichenen Fläche nach der Zeit ist konstant und hängt direkt vom erhaltenen Drehimpuls \(L\) sowie der Masse \(m\) ab, was formal als
\begin{equation}
    \frac{dA}{dt} = \frac{L}{2m}
    \label{kepler:equation:flaechengeschwindigkeit}
\end{equation}
geschrieben wird \cite{astro2012}.
Die Gesamtfläche einer Ellipse berechnet sich aus dem Produkt der Kreiszahl Pi, der grossen Halbachse \(a\) und der kleinen Halbachse \(b\) \cite{astro2012}.
Teilt man diese Gesamtfläche durch die in Gleichung \eqref{kepler:equation:flaechengeschwindigkeit} definierte, konstante Flächengeschwindigkeit, so erhält man die Umlaufzeit \(T\) in der Form
\begin{equation}
    T = \frac{\pi a b}{L / 2m} .
    \label{kepler:equation:umlaufzeit_flaeche}
\end{equation}
Quadriert man diese Gleichung und nutzt die geometrische Eigenschaft \(b^2 = a^2(1-E^2)\) für die Exzentrizität \(E\), lässt sich der Drehimpuls aus der Gleichung eliminieren \cite{astro2012}.
Das finale Resultat entspricht exakt der in Abschnitt \ref{kepler:section:einleitung} gezeigten Proportionalität, welche somit unabhängig von der Exzentrizität der Bahn gültig ist \cite{astro2012}.

